\section{Discussion}
\label{sec:discussion}

The most consistent gains obtained with \rscan{} concern two thermodynamic differences: the balance between solid and liquid Ti, and the enthalpy of an alloy liquid relative to its pure components.
The melting estimate moves toward experiment, and the overall mixing-enthalpy discrepancy decreases, although the broad minimum near equiatomic composition becomes too deep.
Density, phonons, and viscosity do not follow the same ranking.
These observations suggest that the benefit of \rscan{} is associated with how it distinguishes competing atomic environments, rather than with a uniform improvement of all features of the potential-energy surface.

\subsection{The solid--liquid free-energy balance}

At fixed pressure, melting occurs when the Gibbs free energies of the solid and liquid are equal.
Let $g_s^A$ and $g_l^A$ denote these free energies per atom for a reference potential $A$, and let $\delta g_i=g_i^B-g_i^A$ be the change on replacing it by potential $B$.
Linearization about the melting point of $A$ gives
\begin{equation}
 \delta T_m \simeq
 \frac{\delta g_l(T_m^A,p)-\delta g_s(T_m^A,p)}
 {\Delta s_{\mathrm{fus}}^A},
 \label{eq:melting_shift}
\end{equation}
where $\Delta s_{\mathrm{fus}}^A=s_l^A-s_s^A>0$ is the entropy of fusion per atom, evaluated at that point.
Thus, a positive change in the liquid--solid free-energy difference raises the melting temperature.
The approximately 95~K increase from the representative PBE estimate to \rscan{} is consistent with greater relative stability of the crystal near coexistence.
It does not determine whether this change is predominantly energetic or entropic, since both contributions enter $g_l-g_s$.
We have not calculated the fusion enthalpy or entropy separately and therefore do not assign a numerical free-energy correction to this temperature shift.

This distinction is particularly relevant for Ti.
Its high-temperature bcc phase is stabilized by finite-temperature lattice dynamics, which cannot be described by the unstable zero-temperature harmonic spectrum alone.\cite{Souvatzis2008}
Our coexistence calculations start from a bcc slab and sample ionic motion explicitly.
Their melting behavior depends on the full potential-energy surface explored by the two phases, including anharmonic contributions.
The harmonic $\alpha$-Ti spectra characterize a different phase and a local curvature of that surface.
Their mixed ranking is therefore compatible with the clearer improvement of the \rscan{} melting estimate.

The density results likewise show why a simple association between a smaller equilibrium volume and a higher melting point is insufficient.
PBEsol gives the largest Ti densities but the lowest working melting estimate, whereas the PBE and \rscan{} densities at 300~K are nearly identical. The PBE and \rscan{} models also differ appreciably in melting temperature despite much smaller relative differences in liquid density.
Volume probes the pressure derivative of the free energy; melting probes the difference between two free energies.
Errors in these quantities need not vary together.
Related sensitivity of melting predictions to the reference functional has been demonstrated in first-principles calculations accelerated by machine-learned force fields.\cite{Jinnouchi2019}
Here, the remaining underestimation of approximately 88~K by \rscan{} also shows that its improved balance is still imperfect.

\subsection{Mixing energetics and a possible electronic origin}

The mixing enthalpy isolates another difference between environments.
Defining $\delta H_m=H_m^B-H_m^A$ at common temperature and pressure, Eq.~\eqref{eq:mixing} gives
\begin{align}
 \delta\Delta H_{\mathrm{mix}}(x)={}&\delta H_m(x)
 -x\,\delta H_m^{\mathrm{Al}}\nonumber\\
 &-(1-x)\,\delta H_m^{\mathrm{Ti}}.
 \label{eq:mixing_shift}
\end{align}
Any constant energy offsets assigned separately to Ti and Al cancel in this expression.
The change therefore reflects the treatment of alloy environments relative to the elemental liquids, including the different structures sampled by the models.
At $x=0.5$ and 2000~K, replacing PBE by \rscan{} lowers the calculated mixing enthalpy by approximately 6.65~kJ/mol, or 69~meV/atom.
This is a substantial change in the energetic preference for mixing.
It brings much of the composition curve closer to the reproduced calorimetric points, while overshooting their minimum.
The provisional overall improvement consequently supports a more favorable description of this particular energy difference, without implying that increasingly exothermic mixing is always more accurate.

The construction of \rscan{} offers a plausible reason for changes of this kind.
Unlike a GGA, the SCAN family uses the orbital kinetic-energy density in addition to the electron density and its gradient.
This extra ingredient distinguishes electronic environments associated with different bonding regimes.\cite{Sun2015}
The regularization and constraint restoration in \rscan{} improve numerical behavior while retaining much of this flexibility.\cite{Furness2020}
In Ti--Al, changes in chemical coordination and crystal order can therefore receive different exchange-correlation corrections even when average density changes are modest.
This provides a physical basis for expecting shifts in relative energies, but does not prescribe their sign or guarantee improved agreement.

Existing electronic-structure studies illustrate why Ti environments deserve particular attention.
Calculations for crystalline TiAl find a substantial Ti-$3d$ contribution near the Fermi level.\cite{Hussain2011}
Moreover, Song \textit{et al.}\cite{Song1994} found that the alternating Ti and Al planes in $\gamma$-TiAl strengthen particular directional Ti--Ti bonding states relative to hypothetical fcc Ti.
Thus, adding Al can modify the Ti-derived bonding network as well as introduce unlike neighbors.
Changes in Ti--Al hybridization and Ti--Ti bonding are plausible microscopic contributions to the functional dependence of mixing.
The crystalline results motivate this interpretation; they do not establish that the same redistribution persists in a disordered liquid.
Our thermodynamic data contain no direct measurement of orbital occupations, charge transfer, or bond-resolved energies.
An attribution of the \rscan{} improvement to a specific electronic redistribution therefore remains a hypothesis.

The two favorable thermodynamic shifts also cannot be reduced to a single assertion of stronger bonding.
Raising the melting point of pure Ti concerns environments containing only Ti, whereas making alloy mixing more exothermic concerns the alloy relative to both elemental liquids.
Neither shift logically requires the other.
Their coexistence is consistent with improved environment-dependent energy differences, but a common microscopic cause has not been demonstrated.
Broad benchmarks reinforce this qualification: \rscan{} improves many solid-state formation enthalpies, yet PBE remains competitive and can have smaller average errors for intermetallic subsets.\cite{Kothakonda2023}
Those benchmarks concern crystalline formation energies, rather than liquid mixing, but show that the meta-GGA level alone does not establish superiority for a given alloy.

\subsection{Separating reference and representation errors}

The paired phonon calculations make this distinction explicit for the local curvature of the $\alpha$-Ti potential-energy surface.
Direct DFT--experiment comparisons assess the three electronic-structure descriptions without an intervening learned potential.
The corresponding DPA4--DFT comparisons, at identical geometry and phonon settings within each functional, test the fidelity of the final fine-tuned models to those references.
Near $\Gamma$, the closer agreement of PBE DFT with the experimental optical frequencies is a property of the reference calculation, whereas the downward shifts of the nondegenerate optical mode in all three DPA4 models identify an additional model discrepancy.
In particular, the reversed optical-mode ordering for the \rscan{} model is absent from its DFT reference.

For an identifiable mode at a given wave vector and the same calculated geometry, the model--experiment frequency residual is the sum of the model--DFT and DFT--experiment residuals.
The two contributions can reinforce or cancel one another.
Consequently, closer agreement of a learned potential with an experimental point does not necessarily imply more faithful reproduction of its training functional.
The DOS comparison provides a complementary test of that fidelity over the Brillouin zone, showing residual changes in peak positions and spectral weight at common broadening and normalization.

This separates the two levels of validation without making either error source exact.
The DFT--experiment residual includes the harmonic approximation, the use of a static equilibrium cell, and unresolved numerical convergence, in addition to the electronic-structure approximation.
The model--DFT residual tests the final fine-tuned representation under the phonon protocol, which differs from the training-label protocol in electronic sampling and smearing.
It does not isolate the optimizer's contribution or establish the improvement over the pretrained model without a corresponding pretrained phonon calculation.
The remaining optical errors therefore motivate more targeted validation of force derivatives near the crystal, even when pooled energy and force errors are small.

\subsection{Consequences for validation and transfer learning}

The other observables help delimit the thermodynamic interpretation.
At 2000~K, PBE is closer than \rscan{} to the central experimental liquid-Ti density correlation; PBE also gives smaller overall discrepancies for the alloy densities at 1873~K.
More exothermic mixing and stronger volume contraction are therefore not sufficient criteria for choosing the most accurate model.
Likewise, the higher alloy viscosity predicted by \rscan{} does not improve agreement at four of the five experimentally covered compositions shared with the PBE calculations; the Al-rich state at $x_{\mathrm{Al}}=0.90$ is an exception.
For pure Ti, however, \rscan{} is closer to the central experimental viscosity correlation at 2000 and 2100~K, while both models give temperature dependences consistent with smooth Arrhenius behavior over the sampled interval.
The transport comparison therefore depends on composition and temperature as well as on the reference functional.
Viscosity depends on the amplitude and relaxation time of stress fluctuations, with additional sensitivity to density and liquid structure.
A more accurate mixing enthalpy alone cannot determine these dynamical quantities.
Chemical short-range order could help connect energetics and transport, but would require partial pair correlations or explicit chemical-order parameters; it is not established by the enthalpy minimum.

This outcome differs from our Al--Cu study, where changing from PBE to \rscan{} improved the composition dependence of viscosity as well as several thermodynamic properties.\cite{Khazieva2024}
Together, the two systems support testing alternative reference functionals, while showing that a favorable ranking in one alloy does not transfer automatically to another.
Foundation-model fine-tuning makes such comparisons practical by using a common starting representation for the separate DFT datasets.
The present work demonstrates this use of DPA4-Neo, without a separate benchmark of training cost against models trained from scratch.

Beyond the matched harmonic comparison, a quantitative attribution of thermodynamic changes to exchange-correlation physics must still separate functional effects from potential-fitting and sampling errors.
The common-geometry benchmark in Sec.~\ref{sec:results_fit} gives comparable energy- and force-error scales for the three fits, but these pooled errors do not constrain systematic solid--liquid free-energy biases directly.
This cautions against attributing every property difference to the reference functional alone.
Further targeted tests should include common snapshots of hot bcc Ti and the liquid states relevant to coexistence.
Comparisons at common geometry and volume would isolate electronic-energy and stress changes; separate equilibrium ensembles would reveal the contribution of structural relaxation.
Free-energy perturbation or thermodynamic integration, with adequate ensemble overlap, could then test the melting interpretation in Eq.~\eqref{eq:melting_shift}.
Ensemble-averaged projected densities of states and bond-resolved analyses could examine the proposed electronic origin, but would need to be related to these energy differences.

These checks would refine the causal interpretation.
The coexistence and liquid-density calculations show a property-dependent pattern: \rscan{} improves the Ti coexistence melting estimate and the provisional overall mixing-enthalpy comparison, whereas PBE is closer to the available liquid-density references.
The pure-Ti viscosity comparison does not establish a uniform transport ranking.
The finite-size and observation-time limitations of coexistence, the indirect calorimetric reference, and the sampling uncertainty in viscosity set the precision with which this pattern can presently be quantified.
